% 356_SM_Teilchen_GF27_En_ch.tex
% Johann Pascher, 7 September 2026
% Standard Model particles from GF(27)* — overview diagram (landscape A4, Doc. 356)
\providecommand{\BM}{\textbf{[B]}}
\providecommand{\KM}{\textbf{[K]}}
\providecommand{\SM}{\textbf{[S]}}
\providecommand{\XM}{\textbf{[X]}}

\phantomsection
\addcontentsline{toc}{chapter}{%
  Standard Model Particles from GF$(27)^*$ --- Overview Diagram (Doc.~356)}

\section*{Status markers}
\begin{center}
\begin{tabular}{@{}lp{10.5cm}@{}}
\textbf{[B]} & \emph{Proved} --- algebraically established. \\[3pt]
\textbf{[K]} & \emph{Confirmed} --- numerically verified. \\[3pt]
\textbf{[S]} & \emph{Speculative} --- open. \\[3pt]
\textbf{[X]} & \emph{Experimental} --- established; FFGFT derivation open. \\
\end{tabular}
\end{center}

\section*{Structure of the diagram}

Each fermion is displayed as a concentric ring diagram (Doc.~356).
Rings from inside out encode physical properties in order of algebraic depth:

\begin{center}
\begin{tabular}{@{}lp{10cm}@{}}
\textbf{White inner circle} & Particle symbol, charge $Q$, hypercharge $Y$ \\[3pt]
\textbf{Colour ring} (quarks only) & $r/g/b$ sectors = colour charge \BM\ (Doc.~346);
  innermost after white because gluons confine colour (Confinement) \\[3pt]
\textbf{Middle ring} & Left violet $= S_d =$ left-handed \BM;
  right orange $= S_u =$ right-handed \BM\ (Doc.~349);
  grey $= \nu_R$ absent in SM \BM/\SM \\[3pt]
\textbf{Outer ring} & Galois orbit colour: navy = O1/O3 Gen.\,1/3 leptons;
  sky blue = O3 Gen.\,2 leptons; teal = O2 Gen.\,1/3 quarks;
  berry = O4 Gen.\,2 quarks \BM\ (Doc.~348) \\[3pt]
\textbf{Red dashed box} & SU$(2)_L$ doublet pair \BM\ (Doc.~349) \\[3pt]
\textbf{Black dashed circle} & Generation 3 \\
\end{tabular}
\end{center}

\section*{Core result}

All fermion quantum numbers of the Standard Model
($Q$, $Y$, colour charge, chirality, generation structure) follow
algebraically from GF$(27)^*$ without free parameters \BM\
(Docs.~346, 347, 348, 349).

\section{Fermions}

\subsection*{Leptons --- Quadratic Residues (QR)}

Charged leptons ($e$, $\mu$, $\tau$) and neutrinos ($\nu_e$, $\nu_\mu$,
$\nu_\tau$) originate in the QR sector of GF$(27)^*$.
Electric charge $Q \in \{0, -1\}$ follows from the QR bit and
$N \bmod 3$ \BM\ (Doc.~346).
Hypercharge $Y = -1$ follows from the Legendre symbol as NQR bit
(Doc.~347 Theorem~A \BM).
All three generations correspond to the three Frobenius cycles of the
root orbits of order~3 in GF$(27)^*$ \BM\ (Doc.~348 Theorem~A).

\subsection*{Quarks --- Non-Quadratic Residues (NQR)}

Up-type quarks ($u$, $c$, $t$) and down-type quarks ($d$, $s$, $b$)
originate in the NQR sector of GF$(27)^*$.
Electric charges $Q \in \{+2/3, -1/3\}$ follow from the NQR bit
and $N \bmod 3$ without free parameters \BM\ (Doc.~346).
Colour charge $r/g/b$ is forced by the $Z_3$ charge operator
$N \bmod 3$ in GF$(3)$: three colours = three group elements \BM\
(Doc.~346 Theorem~C).

\section{Gauge Bosons}

The 8 gluons arise from $4 \times Z_{13}$-orbits $\times$ $2$
$Z_2$-parities in GF$(27)^*$ \BM\ (Doc.~339) --- the two groups of
four are shown in different shades.
Electroweak symmetry breaking
$SU(2)_L \times U(1)_Y \to U(1)_\text{EM}$ is algebraically forced
by three Galois building blocks \BM\ (Doc.~355).
In FFGFT, $\tilde{T} \cdot m = 1$ replaces the Higgs mechanism
(Doc.~019); the boson is experimentally established \XM.

\section{Open Points}

\begin{itemize}
  \item \textbf{Generation assignment (ordering):} Which orbit element
    corresponds to which generation (1/2/3)? Not derived \SM\ (Doc.~346).
  \item \textbf{CKM/PMNS angles:} Structure established (Doc.~348 Theorems B,C \BM);
    numerical values require further work (Doc.~348 Theorem~D \SM).
  \item \textbf{Higgs derivation:} Mechanism established experimentally \XM;
    formal FFGFT elaboration open (Doc.~019).
\end{itemize}
